\begin{trapbox}
\begin{itemize}
    \item \textbf{量词否定位置错误}：否定"所有……都"时，应改为"存在一个……不"，而不是"所有……都不"。例如"所有人都喜欢数学"的否定是"有人不喜欢数学"，而不是"所有人都不喜欢数学"。
    \item \textbf{比较关系否定遗漏等号}："大于"的否定是"不大于"（即小于等于），学生常误以为"小于"。注意"不大于"包含等于的情况。
    \item \textbf{联结词否定混淆}：$p\land q$ 的否定是 $(\lnot p)\lor(\lnot q)$，容易错写成 $(\lnot p)\land(\lnot q)$。牢记德摩根定律。
    \item \textbf{"至少n个"的否定}：否定是"至多n-1个"，而非"至多n个"或"少于n个"（后者在整数语境下与"至多n-1个"等价，但严格表述应为"至多n-1"）。
    \item \textbf{多重否定}：当原命题本身含有否定词时，需逐层处理，如"并非所有的数都不是整数"等价于"存在一个数是整数"。
\end{itemize}
\end{trapbox}